\chapter{Dizziness}

\section*{Definition and Overview}
Dizziness is a common symptom that people describe in different ways: feeling light-headed, unsteady, faint, or as if the room is spinning. The spinning sensation is called vertigo and usually comes from the inner ear, while light-headedness often relates to blood pressure, heart, or metabolic causes. Most causes of dizziness are benign and short-lived, but dizziness can also signal serious conditions such as heart problems, stroke, or significant blood loss, so sudden or severe dizziness needs prompt assessment. This chapter explains the types of dizziness, their causes, and what to do.

\section*{Epidemiology}
Dizziness is one of the most common symptoms in medical practice, affecting around 20--30\% of the population at some point, and it becomes more common with age. Benign paroxysmal positional vertigo (BPPV), the most common cause of vertigo, affects many older people. Dizziness is a leading reason for GP visits and emergency presentations, and it is a major risk factor for falls in older adults. Many causes are self-limiting, but serious causes must be excluded.

\section*{Aetiology and Pathophysiology}
\begin{itemize}
    \item \textbf{Vertigo:} the sensation of spinning, usually from the inner ear or brain, including BPPV, vestibular neuritis, and Meniere's disease
    \item \textbf{Light-headedness:} feeling faint, often from low blood pressure, dehydration, anaemia, or low blood sugar
    \item \textbf{Cardiac causes:} abnormal heart rhythms, heart valve problems, and heart attack
    \item \textbf{Neurological causes:} stroke, transient ischaemic attack, and migraine
    \item \textbf{Medicines:} many medicines cause dizziness, especially blood pressure and sedative drugs
    \item \textbf{Anxiety:} panic and hyperventilation cause light-headedness
    \item \textbf{Ageing:} reduced balance, vision, and circulation increase dizziness and falls
\end{itemize}

\section*{Clinical Features}
\begin{itemize}
    \item Vertigo: a spinning or rotational sensation, often with nausea and vomiting
    \item Light-headedness: feeling faint or about to pass out, especially on standing
    \item Unsteadiness: feeling off-balance when walking
    \item Associated symptoms: ringing in the ears, hearing loss, headache, palpitations, or chest pain
    \item Symptoms triggered by head movement may indicate BPPV
    \item Dizziness on standing suggests low blood pressure
\end{itemize}

\section*{Differential Diagnosis}
\begin{itemize}
    \item Vertigo from the inner ear versus the brain, which need different management
    \item Benign causes such as BPPV versus serious causes such as stroke
    \item Cardiac arrhythmias and other heart problems
    \item Anaemia, dehydration, and low blood sugar
    \item Medication side effects
    \item Anxiety and panic disorder
\end{itemize}

\section*{When to Seek Medical Care}
See a doctor for recurrent, persistent, or severe dizziness, or dizziness with falls. Seek urgent care for dizziness with chest pain, palpitations, weakness, or when it occurs with a headache or neurological symptoms.

\textbf{Call 000 (or local emergency services) for:}
\begin{itemize}
    \item Dizziness with weakness, facial droop, or difficulty speaking (possible stroke)
    \item Dizziness with chest pain, breathlessness, or palpitations
    \item Fainting, collapse, or loss of consciousness
    \item Dizziness with severe headache, confusion, or a seizure
    \item Dizziness after a head injury
\end{itemize}

\section*{Diagnosis and Investigations}
\begin{itemize}
    \item \textbf{History:} the type of dizziness, triggers, duration, and associated symptoms
    \item \textbf{Physical examination:} blood pressure lying and standing, heart rate, and neurological examination
    \item \textbf{Ear examination and positional testing:} for BPPV and other inner ear causes
    \item \textbf{Blood tests:} blood count, glucose, and electrolytes
    \item \textbf{ECG and heart monitoring:} for arrhythmias
    \item \textbf{Imaging and specialist tests:} for suspected neurological or inner ear causes
\end{itemize}

\section*{Management}
\begin{itemize}
    \item \textbf{BPPV:} treated with the Epley manoeuvre, a series of head movements performed by a clinician
    \item \textbf{Vestibular neuritis:} rest, medicines for nausea, and vestibular rehabilitation
    \item \textbf{Meniere's disease:} low-salt diet, diuretics, and specialist care
    \item \textbf{Low blood pressure:} fluids, salt as advised, and review of medicines
    \item \textbf{Anaemia:} iron treatment and management of the cause
    \item \textbf{Medication review:} reducing or changing medicines that cause dizziness
    \item \textbf{Falls prevention:} balance exercises, home safety, and walking aids
    \item \textbf{Vestibular rehabilitation:} exercises that retrain the balance system
\end{itemize}

\section*{Complications}
\begin{itemize}
    \item Falls, fractures, and head injuries
    \item Loss of independence and confidence in older people
    \item Impaired driving and work performance
    \item Anxiety and avoidance of activity
    \item Missed serious conditions such as stroke or heart disease
\end{itemize}

\section*{Prognosis and Outcomes}
Most causes of dizziness resolve or are well managed. BPPV is readily cured with the Epley manoeuvre in most people, and vestibular neuritis settles over weeks with rehabilitation. Light-headedness from low blood pressure or medicines improves with treatment and adjustment. The main goal is to identify serious causes and prevent falls, and with appropriate care most people regain balance and confidence.

\section*{Prevention and Self-Care}
\begin{itemize}
    \item Stand up slowly from sitting or lying to avoid light-headedness
    \item Stay hydrated and eat regularly
    \item Avoid sudden head movements if prone to vertigo
    \item Review medicines with a doctor if they cause dizziness
    \item Exercise for balance, such as tai chi, and consider vestibular rehabilitation
    \item Make the home safe to prevent falls
    \item Seek prompt care for sudden or severe dizziness
\end{itemize}

\section*{Sources and Further Reading}
The following reputable sources provide further information for healthcare students and patients seeking reliable, current advice about dizziness.

\begin{itemize}
    \item Healthdirect Australia. \textit{Dizziness and vertigo: causes and treatment.} https://www.healthdirect.gov.au/
    \item Vestibular Disorders Association. \textit{Information about balance disorders.} https://vestibular.org/
    \item Balance and Ear Surgery Australia. \textit{Dizziness and balance information.} https://www.balanceear.com.au/
    \item NHS. \textit{Dizziness: causes and when to see a doctor.} https://www.nhs.uk/conditions/dizziness/
    \item Mayo Clinic. \textit{Dizziness: symptoms and causes.} https://www.mayoclinic.org/
    \item MSD Manual. \textit{Dizziness and vertigo: professional overview.} https://www.msdmanuals.com/
\end{itemize}
